\section{Conclusion}
\label{sec:conclusion}

We studied long-horizon search agents by opening the search trajectory rather than treating it as an opaque path to a final answer. Across six agents on BrowseComp-Plus, under a shared ReAct harness and retriever, effort and outcome dissociate: neither search volume nor context budget explains accuracy. What does is whether an agent's queries retrieve the gold evidence, and the behavioral signature of the agents that do is query discipline. The trajectory analysis shows why: useful evidence arrives early or not at all, while later turns run on as a wasted tail of redundant snippets. The failure analysis splits errors into retrieval gaps, where evidence was never surfaced, and utilization gaps, where it was surfaced but not converted into the right answer; a within-question test confirms that this is not merely a capability ranking --- on the same questions, retrieving the gold evidence is what separates success from failure.

These findings argue for making search better directed rather than deeper: for weaker agents the main opportunity is query formulation, recognizing failed hypotheses and avoiding redundant re-querying; for stronger agents it shifts toward verification and answer normalization. Our analysis quantifies the headroom but remains diagnostic: realizing the gains will require training, prompting, or tool-design changes.

% \section{Limitations}
% \label{sec:limitations}

% Our analysis is grounded on a single closed-corpus benchmark, so visit precision coincides numerically with the visit-level redundant rate (the two diverge on the open web). The context-budget composition (Figure~\ref{fig:context-composition}) is sensitive to our $512$-token snippet truncation: shorter open-web snippets ($\sim 30\text{--}50$ tokens) would shift the dominant lever from snippet-stream management toward visit-content management, though the qualitative finding that accumulated retrieval dominates context cost still transfers. The cross-agent design holds the retriever fixed \emph{by design}: a single shared retriever is what attributes the cross-agent recall spread to the agents rather than to retrieval, so we treat the retriever as a controlled constant, not a benchmark target. This leaves two retrieval-side extensions to future work. First, our analysis pins query-side \emph{variance} but not the retriever's absolute recall ceiling; replaying each agent's logged queries under varied retrieval depth and reranking --- holding the queries fixed --- would bound the top-$K{=}10$ contribution (a literal oracle retriever is ill-posed as a query test, since it returns the qrels regardless of the query). Second, whether the query-side findings transfer under a stronger or hybrid embedding model is open: a different retriever would shift the absolute RG/UG levels, though the attribution of cross-agent recall variance to queries should persist as long as the retriever is shared. Document-level relevance and the strict gold-qrels layer $G(q)$ together leave boundary UG cases as partly hidden retrieval gaps, bounded by the boundary UG share in Figure~\ref{fig:failure-combined}c. Finally, we evaluate fixed checkpoints under a unified harness; whether the directional-vs.-last-hop balance shifts under different prompting, tool schemas, or training recipes remains open. Across all these axes the absolute percentages are benchmark- and configuration-specific, but the diagnostic framework itself --- the relevance-grounded RG/UG attribution, the episode-utility decomposition, and the recall-vs-budget contrast --- transfers to any benchmark that provides document-level relevance judgments.
